% Generated draft: requires mathematical and source review.
\documentclass[11pt]{article}
\usepackage[margin=25mm]{geometry}
\usepackage{fontspec}
\setmainfont{Latin Modern Roman}
\usepackage{amsmath,amssymb,mathtools,unicode-math}
\setmathfont{Latin Modern Math}
\usepackage{enumitem,needspace,xurl,hyperref}
\setlength{\parindent}{0pt}
\setlength{\parskip}{5pt}
\setlength{\emergencystretch}{3em}
\newcommand{\N}{\mathbb{N}}
\newcommand{\Z}{\mathbb{Z}}
\newcommand{\Q}{\mathbb{Q}}
\newcommand{\R}{\mathbb{R}}
\newcommand{\C}{\mathbb{C}}
\newcommand{\F}{\mathbb{F}}
\newcommand{\A}{\mathbb{A}}
\newcommand{\PP}{\mathbb{P}}
\newcommand{\E}{\mathbb{E}}
\newcommand{\eps}{\varepsilon}
\newcommand{\dd}{\,\mathrm{d}}
\begin{document}
\section*{UnsolvedMath ID: 20003085}
\textbf{Source problem code:} catalogue entry
\subsection{Definitions and mathematical statement}
Let $S=S_{g,n}$ be an orientable finite-type surface with punctures and $\chi(S)=2-2g-n<0$, and use complete finite-area hyperbolic structures with markings, with no geodesic boundary. This is the explicit setting supplied by the aid; the canonical sentence does not state the surface hypotheses. Let $\mathcal T(S)$ be Teichm\"uller space and let $\mathcal S(S)$ denote essential nonperipheral simple closed curve classes. For $X\in\mathcal T(S)$ and $v\in T_X\mathcal T(S)$, the infinitesimal Thurston asymmetric norm is
\[
 F_X(v)=\sup_{\gamma\in\mathcal S(S)}
 \frac{d_X\ell_\gamma(v)}{\ell_\gamma(X)}
 =\sup_{\gamma\in\mathcal S(S)}d_X\log\ell_\gamma(v).
\]
It is positively homogeneous but need not satisfy $F_X(v)=F_X(-v)$. The zero-dimensional three-punctured-sphere case is interpreted separately as the unique Teichm\"uller point, rather than requiring a nonempty curve supremum.

Dumas and Rafi ask whether the Thurston norm at $X$ determines $X$. The aid interprets this as a Royden-style infinitesimal rigidity question, using surjective \emph{real-linear} norm isometries. It reports that for $X,Y\in\mathcal T(S)$,
\[
 A:T_X\mathcal T(S)\xrightarrow{\cong}T_Y\mathcal T(S),
 \qquad F_Y(Av)=F_X(v)\text{ for every }v,
\]
is precisely a derivative $A=d_Xh$ of an extended mapping class $h\in\operatorname{Map}^{\pm}(S)$ satisfying $hX=Y$. Extended means orientation-reversing classes are allowed; derivatives are understood on the underlying real manifolds. No uniqueness of $h$ is asserted where the mapping-class action has a kernel.

Consequently the reported invariant
\[
 X\longmapsto[(T_X\mathcal T(S),F_X)]_{\mathrm{linear\ isometry}}
\]
has fibers exactly the extended mapping-class orbits, so it determines the underlying \emph{unoriented unmarked} hyperbolic surface. It does not generally recover the marking: distinct marked mapping-class translates have isometric tangent norms. For a zero-dimensional Teichm\"uller space this distinction is vacuous. The statement is about real-linear isometries, not arbitrary nonlinear bijections of tangent sets.

The aid also reports that every real-linear similarity satisfying $F_Y(Av)=cF_X(v)$ with $c>0$ has the form
\[
 A=c\,d_Xh,\qquad hX=Y.
\]
Finally, retaining all curve labels gives a stronger marked conclusion. If a norm isometry $A$ obeys
\[
 A^*(d_Y\log\ell_\gamma)=d_X\log\ell_\gamma
 \qquad\text{for every }\gamma\in\mathcal S(S),
\]
then the aid reports $X=Y$ as marked Teichm\"uller points. The label-preserving condition is additional information on the dual covectors, not part of the abstract norm isometry class. Precise infinitesimal-rigidity and curve-covector theorem hypotheses, including low-complexity cases, must be verified from the missing full references. The registry's solved-in-literature label is an attribution, not an independent verification by this draft.
\subsection{Short English statement}
The supplied answer determines the hyperbolic surface up to an extended mapping class, rather than recovering its marking from an abstract tangent norm. Keeping every curve label is extra data and is reported to recover the marked point. Similarities are stated with their positive scale factor.
\subsection{Sources}
\begin{itemize}
\item Local UnsolvedMath registry, permanent ID 20003085, source code catalogue entry; \texttt{aim-topology-notes.json}, notes[172], zero-based index 172.
\item AIM, \emph{Lipschitz metric on Teichmueller space}, Problem 1.1, attributed to D. Dumas and K. Rafi: \url{https://aimath.org/pastworkshops/lipschitzteichproblems.pdf}.
\item The supplied aid invokes a Royden-style infinitesimal Thurston-rigidity theorem but supplies no author, title or full bibliographic entry for the affirmative result or labeled-covector proposition.
\item Machine-generated research attempt 1: full solution, solved, rigor labelled ``proved,'' low novelty confidence; registry ``SOLVED-IN-LITERATURE.'' Not peer reviewed; independently verify before citation. Registry updated 14 August 2026.
\end{itemize}
\end{document}
